\section{Συνολική Αρχιτεκτονική και Ροή Δεδομένων}
\label{sec:arch_overview}

Το σύστημα αποτελείται από πέντε αυτόνομα υποσυστήματα που επικοινωνούν μέσω αυστηρά καθορισμένων διεπαφών gRPC (Protocol Buffers) και ασύγχρονων ροών μηνυμάτων:
\begin{enumerate}
    \item \textbf{Orchestrators (Rust):} Οδηγοί του εξελικτικού αλγορίθμου, ένας ανά νήσο, που υλοποιούν τον κύκλο $\mu + \lambda$, τη μνήμη γονιδίων και τη διαχείριση του pipeline αξιολόγησης.
    \item \textbf{LEAD DHT Cluster (Rust):} Τριμελής κατανεμημένος δακτύλιος αποθήκευσης και χωρικής ευρετηρίασης γεωμετριών με 60 vnodes.
    \item \textbf{Surrogate Node (Rust + C++):} Υπηρεσία ταχείας πρόβλεψης με ενεργό μοντέλο MLP/GP και ουρά επιγραμμικής επανεκπαίδευσης.
    \item \textbf{Worker Simulator Farm (Python + AeroSandbox):} Συστοιχία υπολογιστικών εργατών που εκτελούν την πλήρη επίλυση 3D VLM.
    \item \textbf{Envoy Load Balancer:} Διαμεσολαβητής αντίστροφης δρομολόγησης gRPC αιτημάτων προς τους εργάτες με αυτόματη δυναμική ανακάλυψη υπηρεσιών.
\end{enumerate}

\subsection*{Λεπτομερής Περιγραφή Υποσυστημάτων και gRPC Συμβολαίων}

Κάθε ένα από τα πέντε υποσυστήματα εκθέτει ένα αυστηρά καθορισμένο gRPC service που ορίζεται μέσω αρχείων \texttt{.proto} και μεταγλωττίζεται αυτόματα κατά την κατασκευή. Ο \textbf{Orchestrator} εκθέτει τη διεπαφή \texttt{OrchestratorService}, η οποία δέχεται αιτήματα \texttt{MigrateElites} από γειτονικούς Orchestrators κατά τη διαδικασία μετανάστευσης, καθώς και \texttt{StatusQuery} για σκοπούς παρακολούθησης. Το \textbf{LEAD DHT Cluster} εκθέτει τη διεπαφή \texttt{LeadStoreService} με μεθόδους \texttt{Put}, \texttt{Get} και \texttt{KNNSearch}: η \texttt{Put} αποθηκεύει ζεύγη γεωμετρίας–απόδοσης, η \texttt{Get} ανακτά γνωστά αποτελέσματα βάσει κλειδιού Hilbert, και η \texttt{KNNSearch} επιστρέφει τους $k$ πλησιέστερους γείτονες ενός ερωτήματος στον ευρετηριασμένο χώρο σχεδίασης. Το \textbf{Surrogate Node} εκθέτει τη διεπαφή \texttt{SurrogateService} με μεθόδους \texttt{Predict} και \texttt{Retrain}: η \texttt{Predict} δέχεται ένα διάνυσμα γονιδίων $\bm{x}^*$ και επιστρέφει τιμή $\hat{y}$ μαζί με αβεβαιότητα $\sigma(\bm{x}^*)$, ενώ η \texttt{Retrain} εκκινεί την ασύγχρονη επανεκπαίδευση. Το \textbf{Worker} εκθέτει τη διεπαφή \texttt{SimulatorService} με μέθοδο \texttt{Evaluate} που λαμβάνει τα αεροδυναμικά παραμετρικά δεδομένα και επιστρέφει τον λόγο $L/D$ μετά από πλήρη επίλυση VLM. Τέλος, ο \textbf{Envoy} δεν ορίζει δικό του \texttt{.proto}, αλλά δρα ως διαφανής πληρεξούσιος (transparent proxy) στρώμα L7, δρομολογώντας κλήσεις προς τη διεπαφή \texttt{SimulatorService} ισορροπώντας το φορτίο μεταξύ όλων των ενεργών εργατών \cite{klein2017envoy}.

\subsection*{Τοπολογία Δικτύου}

Όλα τα υποσυστήματα αναπτύσσονται εντός του κοινού εξωτερικού Docker δικτύου (\texttt{simulated-lan}) και μοιράζονται τον τόμο αρχείων καταγραφής (\texttt{app-logs}), εξασφαλίζοντας απομονωμένη και ρεαλιστική προσομοίωση τοπικού δικτύου (LAN). Ο Envoy εκτίθεται εσωτερικά στη θύρα \textbf{50051} ανά νήσο (και χαρτογραφείται στη θύρα \textbf{50052} του συστήματος για τη δεύτερη νήσο) και αποτελεί το μοναδικό σημείο εισόδου για αιτήματα αξιολόγησης Tier~3 προς τους εργάτες. Οι τρεις κόμβοι του LEAD DHT εκθέτουν HTTP διεπαφές στις θύρες \textbf{2001}, \textbf{2002} και \textbf{2003} και gRPC στις θύρες \textbf{50051}--\textbf{50053}, ενώ ο Surrogate Node ακούει στη θύρα \textbf{50054}. Οι Orchestrators σχηματίζουν ομότιμο δακτύλιο μετανάστευσης στη θύρα \textbf{50060}.

\subsection*{Ροή Δεδομένων κατά την Αξιολόγηση}

Κάθε αίτημα αξιολόγησης γονιδίου ακολουθεί την εξής πορεία: ο \textbf{Orchestrator} κωδικοποιεί το διάνυσμα $\bm{x}^*$ σε μήνυμα Protocol Buffer και υποβάλλει ερώτημα \texttt{KNNSearch} στον τοπικό κόμβο του \textbf{LEAD DHT}, λαμβάνοντας την ελάχιστη απόσταση $d_{\min}$. Εάν το ερώτημα δεν ικανοποιείται από τα Tier~1 ή Tier~2 κριτήρια (βλ.\ Ενότητα~\ref{sec:arch_multi_tier}), αποστέλλεται κλήση gRPC \texttt{Evaluate} στον τοπικό \textbf{Envoy} (θύρα 50051), ο οποίος επιλέγει έναν ελεύθερο \textbf{Worker} και προωθεί το αίτημα. Μόλις ο Worker επιστρέψει αποτέλεσμα, αυτό καταχωρείται στο LEAD DHT μέσω \texttt{Put} και ταυτόχρονα διοχετεύεται στον Surrogate Node (\texttt{Retrain} trigger).

\subsection*{Επιλογή Rust και Τηλεμετρία Kafka}

Η επιλογή της Rust για τα κρίσιμης καθυστέρησης υποσυστήματα (Orchestrator, LEAD DHT, Surrogate) αιτιολογείται από το μοντέλο ιδιοκτησίας μνήμης χωρίς garbage collector, το οποίο εξαλείφει τις διακοπές GC παύσεων που παρατηρούνται σε JVM-βασισμένα συστήματα \cite{matsakis2014rust}. Αυτό εγγυάται προβλέψιμα tail latencies στις κλήσεις gRPC ακόμα και υπό βαρύ φορτίο. Για την παρακολούθηση του συστήματος σε πραγματικό χρόνο, όλα τα εμπορευματοκιβώτια γράφουν δομημένα αρχεία JSON logs στον κοινόχρηστο τόμο \texttt{app-logs}. Ο δαίμονας \textbf{Fluent-Bit} παρακολουθεί τα αρχεία αυτά και δρομολογεί τα συμβάντα βάσει ετικέτας στα αντίστοιχα θέματα (\texttt{logs.lead}, \texttt{logs.orchestrator}, \texttt{logs.worker}, \texttt{logs.surrogate}) ενός \textbf{Kafka} broker σε λειτουργία KRaft \cite{kreps2011kafka}. Στη συνέχεια, ο πίνακας παρακολούθησης \textbf{Monitor Web Dashboard} (Node.js/Express, θύρα 3000) καταναλώνει τα γεγονότα και τα προβάλλει σε πραγματικό χρόνο μέσω WebSockets με διαδραστικό τρισδιάστατο θεατή (3D plane viewer).

Η συνολική αρχιτεκτονική διάταξη του κατανεμημένου συστήματος, οι διασυνδέσεις των υποσυστημάτων και η ροή δεδομένων κατά την εκτέλεση του γενετικού αλγορίθμου απεικονίζονται διαγραμματικά στο Σχήμα~\ref{fig:system_architecture}.

\begin{figure}[htbp]
    \centering
    \includegraphics[width=0.95\textwidth]{figures/fig5_system_architecture.png}
    \caption{Συνολική αρχιτεκτονική τοπολογία του κατανεμημένου συστήματος: Δακτύλιος νήσων Orchestrator (Chord Ring), δακτύλιος LEAD DHT με χωρική ευρετηρίαση Hilbert, επιταχυντής υποκατάστασης Surrogate Node, ισοσταθμιστής φορτίου Envoy και συστοιχία υπολογιστικών εργατών AeroSandbox CFD.}
    \label{fig:system_architecture}
\end{figure}
